\chapter{奇异值分解}

上一讲中我们着重介绍了内积空间上的一些特殊的算子，以及它们的约化分解. 本讲我们将尝试将所得到的结论推广到一般的线性映射上，并给出一个将一般线性映射都“对角化”的方式.

\section{奇异值分解}

我们提出要将一般的线性映射都“对角化”，但在若当标准型的时候我们已经知道，这理应是做不到的，症结在于可能没有足够的特征向量. 所以我们此处的“对角化”并不是使用线性映射本身的特征值和特征向量，而是采取与之相关的特殊的算子的特征值和特征向量. 回忆在对称矩阵部分，我们从一般的矩阵 $A$ 构造出了 $A^{\mathrm{T}}A$ 这样一个对称矩阵. 那么采取类似的方法，我们可以从 $T \in \mathcal{L}(V, W)$ 构造出算子 $T^*T$，让我们来看看它的性质.

\begin{theorem}{$T^*T$ 的性质}{T^*T的性质}
    设 $T \in \mathcal{L}(V, W)$，则
    \begin{enumerate}
        \item $T^*T$ 是正算子.

        \item $\ker T^*T = \ker T$.

        \item $\im T^*T = \im T^*$.

        \item $\dim \im T = \dim \im T^* = \dim \im T^*T$.
    \end{enumerate}
\end{theorem}

\begin{proof}
    \begin{enumerate}
        \item 由伴随的运算性质可知 $(T^*T)^* = T^*(T^*)^* = T^*T$，因此 $T^*T$ 是自伴的. 又对任意 $v \in V$，有 $\langle T^*Tv, v \rangle = \langle Tv, Tv \rangle = \lVert Tv \rVert^2 \geqslant 0$，故 $T^*T$ 是正算子.

        \item 若 $v \in \ker T$，则 $Tv = 0$，从而 $T^*Tv = 0$，故 $\ker T \subseteq \ker T^*T$. 反过来，若 $v \in \ker T^*T$，则 $0 = \langle T^*Tv, v \rangle = \lVert Tv \rVert^2$，因而 $Tv = 0$，即 $v \in \ker T$. 故 $\ker T^*T = \ker T$.

        \item 由\autoref{thm:伴随核空间像空间}可知 $\im T^* = (\ker T)^{\perp}$，又因为 $T^*T$ 是自伴算子，所以 $\im T^*T = (\ker T^*T)^{\perp}$. 结合前述第 2 条，有 $\im T^*T = (\ker T^*T)^{\perp} = (\ker T)^{\perp} = \im T^*$.

        \item 由线性映射基本定理（\autoref{thm:线性映射基本定理}）有 $\dim \im T = \dim V - \dim \ker T$. 同理，结合\autoref{thm:伴随核空间像空间}有 $\dim \im T^* = \dim V - \dim \ker T = \dim \im T$，再由前述第 3 条知 $\dim \im T^*T = \dim \im T^*$，因此 $\dim \im T = \dim \im T^* = \dim \im T^*T$.
    \end{enumerate}
\end{proof}

根据\autoref{thm:T^*T的性质}的第一条，结合正算子特征值非负的性质，我们可以做出如下定义.

\begin{definition}{奇异值}{} \index{qiyizhi@奇异值 (singular value)}
    设 $T \in \mathcal{L}(V, W)$. 则 $T^*T$ 的特征值的算术平方根称为 $T$ 的\term{奇异值}. 我们总将奇异值按降序排列，并且每个奇异值按对应特征子空间的维数重复书写.
\end{definition}

但这样定义出来的数如何描述 $T$ 呢？我们先从线性映射最基本的部分开始.

\begin{theorem}{}{奇异值与单满射}
    设 $T \in \mathcal{L}(V, W)$，则
    \begin{enumerate}
        \item $T$ 是单射当且仅当 $0$ 不是 $T$ 的奇异值.

        \item $T$ 的正奇异值的个数恰为 $\dim \im T$（按重数计算）.

        \item $T$ 是满射当且仅当 $T$ 的正奇异值的个数等于 $\dim W$.
    \end{enumerate}
\end{theorem}

\begin{proof}
    \begin{enumerate}
        \item 由定义知，$0$ 是 $T$ 的奇异值当且仅当 $0$ 是 $T^*T$ 的特征值. 对于正算子而言，这又等价于 $\ker T^*T \neq \{\vec{0}\}$. 由\autoref{thm:T^*T的性质} 第 2 条可知
              \[
                  \ker T^*T = \ker T,
              \]
              故 $0$ 不是 $T$ 的奇异值当且仅当 $\ker T = \{\vec{0}\}$，也即 $T$ 是单射.

        \item 设 $\dim V = n$. 则 $T^*T$ 的 $0$ 特征值的重数为
              \[
                  \dim \ker T^*T = \dim \ker T.
              \]
              因而正奇异值的个数为
              \[
                  n - \dim \ker T = \dim \im T,
              \]
              最后一步由线性映射基本定理得到.

        \item $T$ 满射当且仅当 $\dim \im T = \dim W$. 再结合第 2 条，便得到结论.
    \end{enumerate}
\end{proof}

特征值对应有特征向量，奇异值也对应有奇异向量.

\begin{definition}{奇异向量}{} \index{qiyixiangliang@奇异向量 (singular vector)}
    设 $T \in \mathcal{L}(V, W)$，$s$ 是 $T$ 的一个奇异值. 若 $v \in V, w \in W$ 满足
    \[
        \lVert v \rVert = \lVert w \rVert = 1,\qquad Tv = sw,\qquad T^*w = sv,
    \]
    则称 $v$ 为 $T$ 关于 $s$ 的\term{右奇异向量}\index{youqiyixiangliang@右奇异向量 (right singular vector)}，称 $w$ 为 $T$ 关于 $s$ 的\term{左奇异向量}\index{zuoqiyixiangliang@左奇异向量 (left singular vector)}.
\end{definition}

$T^*T$ 除去其特征值的性质外，更重要的在于其是一个自伴算子，所以我们可以对它应用谱定理，并由此给出我们所谓的“对角化”的方式——奇异值分解.

\begin{theorem}{奇异值分解}{奇异值分解定理} \index{qiyizhi!fenjie@奇异值分解 (singular value decomposition)}
    设 $T \in \mathcal{L}(V, W)$，$\dim V = n,\enspace \dim W = p$，其正奇异值按重数计为
    \[
        s_1 \geqslant s_2 \geqslant \cdots \geqslant s_r > 0.
    \]
    则存在 $V$ 的一组标准正交基 $e_1, \ldots, e_n$ 和 $W$ 的一组标准正交基 $f_1, \ldots, f_p$，使得
    \[
        Te_k = s_k f_k,\enspace k = 1, \ldots, r;\qquad Te_k = 0,\enspace k = r+1, \ldots, n.
    \]
    因而对任意 $v \in V$ 都有
    \[
        Tv = s_1 \langle v, e_1 \rangle f_1 + \cdots + s_r \langle v, e_r \rangle f_r.
    \]
\end{theorem}

\begin{proof}
    因为 $T^*T$ 是正算子，所以它是自伴的. 由谱定理，存在 $V$ 的一组标准正交基 $e_1, \ldots, e_n$，使得
    \[
        T^*Te_k = s_k^2 e_k,\enspace k = 1, \ldots, n,
    \]
    其中约定 $s_{r+1} = \cdots = s_n = 0$.

    对每个 $k = 1, \ldots, r$，定义
    \[
        f_k = \frac{1}{s_k}Te_k.
    \]
    下证 $f_1, \ldots, f_r$ 是 $W$ 中的标准正交组. 对任意 $j, k \in \{1, \ldots, r\}$，有
    \begin{align*}
        \langle f_j, f_k \rangle
         & = \frac{1}{s_j s_k}\langle Te_j, Te_k \rangle
          = \frac{1}{s_j s_k}\langle e_j, T^*Te_k \rangle \\
         & = \frac{1}{s_j s_k}\langle e_j, s_k^2 e_k \rangle
          = \frac{s_k}{s_j}\langle e_j, e_k \rangle
          = \delta_{jk}.
    \end{align*}
    因此 $f_1, \ldots, f_r$ 的确是 $W$ 的标准正交组.

    对于 $k = r+1, \ldots, n$，因为 $T^*Te_k = 0$，由\autoref{thm:T^*T的性质} 第 2 条知 $e_k \in \ker T$，从而
    \[
        Te_k = 0.
    \]

    接下来将 $f_1, \ldots, f_r$ 扩充为 $W$ 的一组标准正交基 $f_1, \ldots, f_p$. 对任意 $v \in V$，由于 $e_1, \ldots, e_n$ 是 $V$ 的一组标准正交基，故
    \[
        v = \langle v, e_1 \rangle e_1 + \cdots + \langle v, e_n \rangle e_n.
    \]
    于是
    \begin{align*}
        Tv
         & = \langle v, e_1 \rangle Te_1 + \cdots + \langle v, e_n \rangle Te_n \\
         & = s_1 \langle v, e_1 \rangle f_1 + \cdots + s_r \langle v, e_r \rangle f_r.
    \end{align*}
    命题得证.
\end{proof}

这一结果告诉我们：对于一般线性映射而言，虽然未必能在同一个空间里找到一组基把它化成对角矩阵，但如果允许在定义域和值域分别选取不同的标准正交基，那么它依然可以被化成极其简单的形式. 这正是奇异值分解与若当标准型并不矛盾的原因.

由此我们也可以得出 $T$ 在这两组标准正交基下的矩阵表示. 若将 $T$ 在 $(e_1, \ldots, e_n)$ 与 $(f_1, \ldots, f_p)$ 下的矩阵记为 $\Sigma = (\sigma_{ij})_{p \times n}$，则
\[
    \sigma_{ij} =
    \begin{cases}
        s_i, & 1 \leqslant i = j \leqslant r, \\
        0,   & \text{其他情况}.
    \end{cases}
\]
它的形式与对角矩阵十分类似，只是可能不是方阵，我们将其定义为矩形对角矩阵.

\begin{definition}{矩形对角矩阵}{} \index{juxingduijiaojuzhen@矩形对角矩阵 (rectangular diagonal matrix)}
    设 $A = (a_{ij})_{m \times n}$，若当 $i \neq j$ 时总有 $a_{ij} = 0$，则称 $A$ 为\term{矩形对角矩阵}.
\end{definition}

通过线性映射的矩阵表示过程，我们也可以得到矩阵形式的奇异值分解.

\begin{theorem}{矩阵的奇异值分解}{矩阵的奇异值分解}
    设 $A$ 为一 $m \times n$ 阶复矩阵，则存在 $m \times m$ 阶酉矩阵 $U$、$n \times n$ 阶酉矩阵 $V$ 以及 $m \times n$ 阶非负实数矩形对角矩阵 $\Sigma$，使得
    \[
        A = U \Sigma V^*.
    \]
    这称为矩阵 $A$ 的奇异值分解. 特别地，当 $A$ 是实矩阵时，$U, V$ 可以取为正交矩阵，此时分解写成
    \[
        A = U \Sigma V^{\mathrm{T}}.
    \]
\end{theorem}

\begin{proof}
    将 $A$ 视为标准内积空间之间的线性映射
    \[
        T_A : \mathbf{C}^n \to \mathbf{C}^m,\qquad T_A(x) = Ax.
    \]
    由\autoref{thm:奇异值分解定理}，存在 $\mathbf{C}^n$ 的一组标准正交基 $v_1, \ldots, v_n$ 和 $\mathbf{C}^m$ 的一组标准正交基 $u_1, \ldots, u_m$，使得
    \[
        Av_k = s_k u_k,\enspace k = 1, \ldots, r;\qquad Av_k = 0,\enspace k = r+1, \ldots, n.
    \]
    令
    \[
        U = (u_1, \ldots, u_m),\qquad V = (v_1, \ldots, v_n).
    \]
    因为两组列向量都是标准正交基，所以 $U, V$ 分别是酉矩阵（实数域上就是正交矩阵）.

    设 $\Sigma$ 为 $m \times n$ 阶矩形对角矩阵，其对角线上前 $r$ 个元素依次为 $s_1, \ldots, s_r$，其余元素全为 $0$. 则
    \[
        AV = (Av_1, \ldots, Av_n) = (s_1u_1, \ldots, s_ru_r, 0, \ldots, 0) = U\Sigma.
    \]
    两边右乘 $V^*$ 即得
    \[
        A = U\Sigma V^*.
    \]
    命题得证.
\end{proof}

这也表明所有矩阵的作用都可以看成三步：先在定义域中做一次保长度的变换 $V^*$，再沿着若干彼此正交的方向做伸缩 $\Sigma$，最后在值域中做一次保长度的变换 $U$.

\begin{example}{求矩阵的奇异值分解}{奇异值分解步骤例题}
    求矩阵
    \[
        A = \begin{pmatrix}
            1 & 1  \\
            1 & 1  \\
            1 & -1
        \end{pmatrix}
    \]
    的奇异值分解.
\end{example}

\begin{solution}
    我们先总结一套求矩阵奇异值分解的通用步骤，然后在本题中逐步执行.

    \begin{enumerate}
        \item 先计算 $A^{\mathrm{T}}A$（或复数域中的 $A^*A$），求出其全部特征值与一组标准正交特征向量.

        \item 将这些特征值开平方，得到奇异值，并按从大到小排序.

        \item 与各个正奇异值 $s_i$ 对应的单位特征向量 $v_i$ 就是右奇异向量，再由公式
              \[
                  u_i = \frac{1}{s_i}Av_i
              \]
              求出左奇异向量.

        \item 若 $U$ 或 $V$ 的列向量还不够构成一组标准正交基，就将其补齐.

        \item 以左奇异向量为列组成 $U$，以右奇异向量为列组成 $V$，再把奇异值按顺序填入 $\Sigma$ 的主对角线，即得
              \[
                  A = U\Sigma V^{\mathrm{T}}.
              \]
    \end{enumerate}

    下面按此步骤计算.

    \begin{enumerate}
        \item 先求 $A^{\mathrm{T}}A$：
              \[
                  A^{\mathrm{T}}A =
                  \begin{pmatrix}
                      1 & 1 & 1  \\
                      1 & 1 & -1
                  \end{pmatrix}
                  \begin{pmatrix}
                      1 & 1  \\
                      1 & 1  \\
                      1 & -1
                  \end{pmatrix}
                  =
                  \begin{pmatrix}
                      3 & 1 \\
                      1 & 3
                  \end{pmatrix}.
              \]

              其特征多项式为
              \[
                  \left| \lambda E - A^{\mathrm{T}}A \right|
                  =
                  \begin{vmatrix}
                      \lambda - 3 & -1          \\
                      -1          & \lambda - 3
                  \end{vmatrix}
                  =
                  (\lambda - 3)^2 - 1
                  =
                  (\lambda - 4)(\lambda - 2).
              \]
              所以特征值为 $4$ 与 $2$.

              对应地，
              \[
                  \lambda = 4 \text{ 时，可取特征向量 } (1, 1)^{\mathrm{T}},
              \]
              \[
                  \lambda = 2 \text{ 时，可取特征向量 } (1, -1)^{\mathrm{T}}.
              \]
              归一化后得到一组标准正交特征向量
              \[
                  v_1 = \frac{1}{\sqrt{2}}\begin{pmatrix}
                      1 \\
                      1
                  \end{pmatrix},\qquad
                  v_2 = \frac{1}{\sqrt{2}}\begin{pmatrix}
                      1  \\
                      -1
                  \end{pmatrix}.
              \]

        \item 因此 $A$ 的奇异值为
              \[
                  s_1 = \sqrt{4} = 2,\qquad s_2 = \sqrt{2}.
              \]

        \item 由 $u_i = \dfrac{1}{s_i}Av_i$ 求左奇异向量.

              对于 $v_1$，
              \[
                  Av_1
                  =
                  \frac{1}{\sqrt{2}}
                  \begin{pmatrix}
                      1 & 1  \\
                      1 & 1  \\
                      1 & -1
                  \end{pmatrix}
                  \begin{pmatrix}
                      1 \\
                      1
                  \end{pmatrix}
                  =
                  \frac{1}{\sqrt{2}}
                  \begin{pmatrix}
                      2 \\
                      2 \\
                      0
                  \end{pmatrix}
                  =
                  \begin{pmatrix}
                      \sqrt{2} \\
                      \sqrt{2} \\
                      0
                  \end{pmatrix},
              \]
              故
              \[
                  u_1 = \frac{1}{2}Av_1
                  =
                  \frac{1}{\sqrt{2}}
                  \begin{pmatrix}
                      1 \\
                      1 \\
                      0
                  \end{pmatrix}.
              \]

              对于 $v_2$，
              \[
                  Av_2
                  =
                  \frac{1}{\sqrt{2}}
                  \begin{pmatrix}
                      1 & 1  \\
                      1 & 1  \\
                      1 & -1
                  \end{pmatrix}
                  \begin{pmatrix}
                      1  \\
                      -1
                  \end{pmatrix}
                  =
                  \frac{1}{\sqrt{2}}
                  \begin{pmatrix}
                      0 \\
                      0 \\
                      2
                  \end{pmatrix}
                  =
                  \begin{pmatrix}
                      0      \\
                      0      \\
                      \sqrt{2}
                  \end{pmatrix},
              \]
              故
              \[
                  u_2 = \frac{1}{\sqrt{2}}Av_2
                  =
                  \begin{pmatrix}
                      0 \\
                      0 \\
                      1
                  \end{pmatrix}.
              \]

        \item 由于 $A$ 是 $3 \times 2$ 矩阵，$U$ 还需要第三个列向量. 我们只需补一个与 $u_1, u_2$ 都正交的单位向量即可，例如取
              \[
                  u_3 = \frac{1}{\sqrt{2}}
                  \begin{pmatrix}
                      1  \\
                      -1 \\
                      0
                  \end{pmatrix}.
              \]
              于是
              \[
                  U =
                  \begin{pmatrix}
                      \frac{1}{\sqrt{2}} & 0 & \frac{1}{\sqrt{2}}  \\
                      \frac{1}{\sqrt{2}} & 0 & -\frac{1}{\sqrt{2}} \\
                      0                  & 1 & 0
                  \end{pmatrix},\qquad
                  V =
                  \frac{1}{\sqrt{2}}
                  \begin{pmatrix}
                      1 & 1  \\
                      1 & -1
                  \end{pmatrix}.
              \]

        \item 最后构造
              \[
                  \Sigma =
                  \begin{pmatrix}
                      2 & 0        \\
                      0 & \sqrt{2} \\
                      0 & 0
                  \end{pmatrix}.
              \]
              因而得到
              \[
                  A = U\Sigma V^{\mathrm{T}}.
              \]
    \end{enumerate}

    至此，奇异值分解求得.
\end{solution}

\section{奇异值分解的应用}

\subsection{奇异值的几何解释}

奇异值分解最重要的意义之一，在于它把一个一般的线性变换分解成了“旋转（或镜像）$+$ 拉伸 $+$ 旋转（或镜像）”三步. 因而奇异值并不是某种纯粹代数的数字，而是直接描述线性变换几何行为的量.

为了叙述方便，本小节只讨论实矩阵的情形. 设
\[
    A = U\Sigma V^{\mathrm{T}} \in \mathbf{R}^{m \times n},
\]
其中奇异值为
\[
    s_1 \geqslant s_2 \geqslant \cdots \geqslant s_r > 0.
\]
记 $B = \{x \in \mathbf{R}^n \mid \lVert x \rVert \leqslant 1\}$ 为 $\mathbf{R}^n$ 的单位球. 我们来研究 $A(B)$ 的形状.

\begin{theorem}{}{奇异值的椭球解释}
    设 $A = U\Sigma V^{\mathrm{T}} \in \mathbf{R}^{m \times n}$ 的正奇异值为 $s_1, \ldots, s_r$. 则单位球 $B$ 在 $A$ 下的像是位于 $\im A$ 中的一个（可能退化的）椭球：
    \[
        A(B) =
        \left\{
            U
            \begin{pmatrix}
                y_1 \\
                \vdots \\
                y_m
            \end{pmatrix}
            \,\middle|\,
            \frac{y_1^2}{s_1^2} + \cdots + \frac{y_r^2}{s_r^2} \leqslant 1,\enspace
            y_{r+1} = \cdots = y_m = 0
        \right\}.
    \]
    也就是说，奇异值 $s_1, \ldots, s_r$ 正好就是这个椭球各条半轴的长度，而左奇异向量给出这些半轴在值域中的方向.
\end{theorem}

\begin{proof}
    任取 $x \in B$，令 $z = V^{\mathrm{T}}x$. 由于 $V$ 是正交矩阵，所以
    \[
        \lVert z \rVert = \lVert x \rVert \leqslant 1.
    \]
    又因为
    \[
        Ax = U\Sigma V^{\mathrm{T}}x = U\Sigma z,
    \]
    所以只需研究 $\Sigma$ 对单位球的作用.

    设
    \[
        z = (z_1, \ldots, z_n)^{\mathrm{T}},
    \]
    则
    \[
        \Sigma z = (s_1z_1, \ldots, s_rz_r, 0, \ldots, 0)^{\mathrm{T}}.
    \]
    若记
    \[
        y_i = s_i z_i,\enspace i = 1, \ldots, r,
    \]
    则由 $z_1^2 + \cdots + z_n^2 \leqslant 1$ 可得
    \[
        \frac{y_1^2}{s_1^2} + \cdots + \frac{y_r^2}{s_r^2}
        = z_1^2 + \cdots + z_r^2
        \leqslant 1.
    \]
    反过来，若 $y_1, \ldots, y_r$ 满足
    \[
        \frac{y_1^2}{s_1^2} + \cdots + \frac{y_r^2}{s_r^2} \leqslant 1,
    \]
    则取
    \[
        z = \left( \frac{y_1}{s_1}, \ldots, \frac{y_r}{s_r}, 0, \ldots, 0 \right)^{\mathrm{T}},
    \]
    便有 $\lVert z \rVert \leqslant 1$，从而对应某个 $x = Vz \in B$.

    这就说明 $\Sigma(B)$ 正是一个坐标轴方向上的椭球，而左乘正交矩阵 $U$ 只是把这个椭球整体旋转（或镜像）到值域中的另一个位置，不改变半轴长度. 因此 $A(B)$ 是一个半轴长度分别为 $s_1, \ldots, s_r$ 的椭球，且这些半轴方向恰好由 $U$ 的前 $r$ 个列向量，也就是左奇异向量给出.
\end{proof}

这一结论十分值得反复体会：
\begin{enumerate}
    \item 右奇异向量给出定义域中彼此正交的“特殊方向”.

    \item 线性变换 $A$ 在这些方向上的伸张倍数分别就是奇异值.

    \item 左奇异向量给出这些方向在值域中的落点方向.

    \item 最大奇异值对应最强的拉伸方向，最小的正奇异值对应最弱的非零拉伸方向.
\end{enumerate}

下面这张示意图展示了二维情形下“单位圆 $\to$ 轴对齐椭圆 $\to$ 旋转后的椭圆”的过程.

\begin{center}
    \begin{tikzpicture}[scale=1.0, >=latex]
        \tikzset{every path/.style={thick}}

        \begin{scope}[xshift=-5.2cm]
            \draw[->] (-1.6,0) -- (1.8,0) node[right] {$x_1$};
            \draw[->] (0,-1.4) -- (0,1.6) node[above] {$x_2$};
            \draw (0,0) circle (1);
            \draw[->,blue] (0,0) -- (0.82,0.57) node[above right] {$v_1$};
            \draw[->,blue] (0,0) -- (-0.57,0.82) node[above left] {$v_2$};
            \node at (0,-1.8) {单位圆};
        \end{scope}

        \begin{scope}
            \draw[->] (-2.2,0) -- (2.4,0);
            \draw[->] (0,-1.4) -- (0,1.6);
            \draw (0,0) ellipse (1.7 and 0.8);
            \draw[dashed] (-1.7,0) -- (1.7,0);
            \draw[dashed] (0,-0.8) -- (0,0.8);
            \node[above right] at (1.7,0) {$s_1$};
            \node[right] at (0,0.8) {$s_2$};
            \node at (0,-1.8) {经 $\Sigma$ 拉伸后};
        \end{scope}

        \begin{scope}[xshift=5.4cm]
            \draw[->] (-2.1,0) -- (2.3,0);
            \draw[->] (0,-1.4) -- (0,1.6);
            \draw[rotate=28] (0,0) ellipse (1.7 and 0.8);
            \draw[->,red] (0,0) -- ({1.5*cos(28)},{1.5*sin(28)}) node[above right] {$u_1$};
            \draw[->,red] (0,0) -- ({-0.7*sin(28)},{0.7*cos(28)}) node[above left] {$u_2$};
            \node at (0,-1.8) {再经 $U$ 旋转后};
        \end{scope}
    \end{tikzpicture}
\end{center}

在这个意义下，奇异值分解把一个一般矩阵的几何本质完全揭示了出来：它无非就是先在定义域中调整观察方向，再沿着若干互相正交的方向做不同比例的拉伸，最后再把结果旋转到值域中的适当位置.

\subsection{极分解定理}

我们继续算子与数的类比. 每个非零复数 $z$ 都可以写成如下形式
\[
    z = \left(\frac{z}{\lvert z \rvert}\right)\lvert z \rvert
      = \left(\frac{z}{\lvert z \rvert}\right)\sqrt{\overline{z}z},
\]
其中 $\dfrac{z}{\lvert z \rvert}$ 是一个单位复数. 那么由此我们可以类比得出一个对 $V$ 上任意算子的漂亮的描述，因其形式类似于复数极坐标分解，所以称为\term{极分解定理}.

\begin{theorem}{极分解定理}{极分解定理} \index{jifenjiedingli@极分解定理 (polar decomposition theorem)}
    设 $T \in \mathcal{L}(V)$，则存在一个等距同构 $S \in \mathcal{L}(V)$，使得
    \[
        T = S\sqrt{T^*T}.
    \]
\end{theorem}

\begin{proof}
    记
    \[
        R = \sqrt{T^*T}.
    \]
    由于 $T^*T$ 是正算子，所以 $R$ 的定义是合理的.

    对任意 $v \in V$，有
    \begin{align*}
        \lVert Tv \rVert^2
         & = \langle Tv, Tv \rangle
          = \langle T^*Tv, v \rangle
          = \langle R^2v, v \rangle  \\
         & = \langle Rv, Rv \rangle
          = \lVert Rv \rVert^2.
    \end{align*}
    因此
    \[
        \lVert Tv \rVert = \lVert Rv \rVert,\enspace \forall v \in V.
    \]

    由此定义映射
    \[
        S_1 : \im R \to \im T,\qquad S_1(Rv) = Tv.
    \]
    先验证 $S_1$ 是良定义的. 若 $Rv_1 = Rv_2$，则
    \[
        R(v_1-v_2)=0.
    \]
    从而
    \[
        \lVert T(v_1-v_2) \rVert = \lVert R(v_1-v_2) \rVert = 0,
    \]
    故 $Tv_1 = Tv_2$，所以 $S_1$ 良定义.

    又对任意 $u = Rv \in \im R$，有
    \[
        \lVert S_1u \rVert = \lVert Tv \rVert = \lVert Rv \rVert = \lVert u \rVert.
    \]
    因而 $S_1$ 保持范数，特别地它是单射. 另一方面，由定义可知 $S_1$ 满射到 $\im T$，所以
    \[
        \dim \im R = \dim \im T.
    \]

    于是
    \[
        \dim (\im R)^\perp = \dim (\im T)^\perp.
    \]
    取 $(\im R)^\perp$ 的一组标准正交基 $e_1, \ldots, e_m$ 和 $(\im T)^\perp$ 的一组标准正交基 $f_1, \ldots, f_m$. 定义
    \[
        S_2 : (\im R)^\perp \to (\im T)^\perp,\qquad
        S_2(a_1e_1 + \cdots + a_me_m) = a_1f_1 + \cdots + a_mf_m.
    \]
    由定义可知，$S_2$ 将一组标准正交基映为另一组标准正交基，因此它也是保持范数的线性同构.

    由于
    \[
        V = \im R \oplus (\im R)^\perp,
    \]
    所以任意 $v \in V$ 可唯一写为
    \[
        v = u + w,\qquad u \in \im R,\enspace w \in (\im R)^\perp.
    \]
    定义
    \[
        Sv = S_1u + S_2w.
    \]
    则 $S \in \mathcal{L}(V)$. 此时对任意 $v \in V$，
    \[
        S(Rv) = S_1(Rv) = Tv,
    \]
    因而
    \[
        T = SR.
    \]

    最后证明 $S$ 是等距同构. 因为 $S_1u \in \im T$，$S_2w \in (\im T)^\perp$，二者正交，所以
    \begin{align*}
        \lVert Sv \rVert^2
         & = \lVert S_1u + S_2w \rVert^2 \\
         & = \lVert S_1u \rVert^2 + \lVert S_2w \rVert^2 \\
         & = \lVert u \rVert^2 + \lVert w \rVert^2 \\
         & = \lVert v \rVert^2.
    \end{align*}
    因此 $S$ 保持范数，从而保持内积，是一个等距同构. 命题得证.
\end{proof}

极分解定理将所有一般的算子表成了一个等距同构和一个正算子的乘积，而这两种算子我们之前的章节都已经给出了完全的描述. 所以我们也得到了一般算子的第二种刻画方式.

\subsection{特征人脸算法}

奇异值分解在经典的人脸识别算法——特征人脸算法（eigenfaces）中有着重要应用. 这一算法的核心思想并不神秘：它本质上就是将高维图像数据投影到一个维数较低、但尽可能保留主要信息的子空间中，而这个子空间正是由主成分分析（PCA）给出的.

为了叙述清楚，我们引入如下记号. 设一共有 $m$ 张人脸图像，每张图像拉直后看作一个 $n$ 维列向量
\[
    x_1, x_2, \ldots, x_m \in \mathbf{R}^n.
\]
通常 $n$ 很大，因为它等于图像的像素数. 将这些样本按行排成矩阵
\[
    D =
    \begin{pmatrix}
        x_1^{\mathrm{T}} \\
        \vdots           \\
        x_m^{\mathrm{T}}
    \end{pmatrix}
    \in \mathbf{R}^{m \times n}.
\]
记样本均值为
\[
    \overline{x} = \frac{1}{m}\sum_{i=1}^m x_i,
\]
并定义中心化后的数据矩阵
\[
    A =
    \begin{pmatrix}
        (x_1-\overline{x})^{\mathrm{T}} \\
        \vdots                          \\
        (x_m-\overline{x})^{\mathrm{T}}
    \end{pmatrix}
    \in \mathbf{R}^{m \times n}.
\]

在这种记号下，样本协方差矩阵为
\[
    \mathrm{Cov}(X)
    = \frac{1}{m-1}\sum_{i=1}^m (x_i-\overline{x})(x_i-\overline{x})^{\mathrm{T}}
    = \frac{1}{m-1}A^{\mathrm{T}}A.
\]
主成分分析告诉我们：如果希望用一个 $k$ 维子空间去近似表示这些图像，并使投影误差尽可能小，那么应当选取协方差矩阵 $A^{\mathrm{T}}A$ 的前 $k$ 个最大特征值所对应的单位特征向量，作为这个子空间的一组标准正交基. 在人脸识别中，这些特征向量所对应的“基脸”就称为特征人脸.

问题在于：在实际应用中，图像张数 $m$ 通常只有几百，而像素数 $n$ 往往达到几万. 因此直接对 $n \times n$ 的矩阵 $A^{\mathrm{T}}A$ 做特征值分解，计算代价非常高. 这时奇异值分解就发挥作用了.

设
\[
    A = U\Sigma V^{\mathrm{T}}
\]
是 $A$ 的奇异值分解，其中
\[
    U \in \mathbf{R}^{m \times m},\qquad V \in \mathbf{R}^{n \times n}
\]
为正交矩阵，$\Sigma \in \mathbf{R}^{m \times n}$ 为非负实数矩形对角矩阵. 则
\[
    AA^{\mathrm{T}} = U(\Sigma \Sigma^{\mathrm{T}})U^{\mathrm{T}},\qquad
    A^{\mathrm{T}}A = V(\Sigma^{\mathrm{T}}\Sigma)V^{\mathrm{T}}.
\]
这说明：
\begin{enumerate}
    \item $U$ 的列向量是 $AA^{\mathrm{T}}$ 的标准正交特征向量.

    \item $V$ 的列向量是 $A^{\mathrm{T}}A$ 的标准正交特征向量.

    \item $A^{\mathrm{T}}A$ 与 $AA^{\mathrm{T}}$ 具有相同的正特征值，这些正特征值恰好是奇异值的平方.
\end{enumerate}

因此，当 $m \ll n$ 时，我们完全可以先求较小的 $m \times m$ 矩阵 $AA^{\mathrm{T}}$ 的特征值分解. 设
\[
    AA^{\mathrm{T}}u_i = \lambda_i u_i,\qquad \lambda_i > 0,
\]
其中 $u_i$ 已单位化. 对应的奇异值为
\[
    s_i = \sqrt{\lambda_i},
\]
并且可以由
\[
    v_i = \frac{1}{s_i}A^{\mathrm{T}}u_i
\]
恢复出 $A^{\mathrm{T}}A$ 的单位特征向量 $v_i$. 这些 $v_i$ 就是我们真正需要的特征人脸方向.

整理成算法步骤，特征人脸方法可以概括为：
\begin{enumerate}
    \item 将每张人脸图像写成向量，并做中心化.

    \item 构造矩阵 $A$ 与较小的矩阵 $AA^{\mathrm{T}}$.

    \item 计算 $AA^{\mathrm{T}}$ 的前 $k$ 个最大特征值及其单位特征向量 $u_1, \ldots, u_k$.

    \item 通过
          \[
              v_i = \frac{1}{\sqrt{\lambda_i}}A^{\mathrm{T}}u_i
          \]
          恢复出特征人脸 $v_1, \ldots, v_k$.

    \item 将任意新的人脸样本 $x$ 投影到子空间
          \[
              \spa(v_1, \ldots, v_k)
          \]
          中，用投影系数作为低维特征用于识别、分类或压缩.
\end{enumerate}

从线性代数的角度看，特征人脸算法之所以高效，正是因为奇异值分解把“求大矩阵 $A^{\mathrm{T}}A$ 的主特征向量”这一任务，转化成了“求小矩阵 $AA^{\mathrm{T}}$ 的主特征向量再回代恢复”的任务. 这也是奇异值分解在数据分析中极其重要的原因之一.

\subsection*{算子范数}

奇异值除了刻画各个特殊方向上的伸缩倍数之外，还控制了线性映射整体上“能够把向量拉多长”. 这就引出算子范数的概念.

\begin{definition}{算子范数}{} \index{suanzifanshu@算子范数 (operator norm)}
    设 $T \in \mathcal{L}(V, W)$，定义
    \[
        \lVert T \rVert = \max_{\lVert v \rVert = 1}\lVert Tv \rVert.
    \]
    称 $\lVert T \rVert$ 为 $T$ 的\term{算子范数}.
\end{definition}

\begin{theorem}{}{}
    设 $T \in \mathcal{L}(V, W)$ 的奇异值按重数记为
    \[
        s_1 \geqslant s_2 \geqslant \cdots \geqslant s_n \geqslant 0,
    \]
    其中 $n = \dim V$. 则
    \[
        \lVert T \rVert = s_1.
    \]
\end{theorem}

\begin{proof}
    由\autoref{thm:奇异值分解定理}，存在 $V$ 的一组标准正交基 $e_1, \ldots, e_n$ 与 $W$ 的标准正交组 $f_1, \ldots, f_n$，使得对任意 $v \in V$，
    \[
        Tv = s_1 \langle v, e_1 \rangle f_1 + \cdots + s_n \langle v, e_n \rangle f_n.
    \]

    任取满足 $\lVert v \rVert = 1$ 的向量，写成
    \[
        v = a_1e_1 + \cdots + a_ne_n.
    \]
    因为 $e_1, \ldots, e_n$ 是标准正交基，所以
    \[
        |a_1|^2 + \cdots + |a_n|^2 = 1.
    \]
    从而
    \[
        Tv = s_1a_1f_1 + \cdots + s_na_nf_n,
    \]
    于是
    \[
        \lVert Tv \rVert^2 = s_1^2|a_1|^2 + \cdots + s_n^2|a_n|^2 \leqslant s_1^2(|a_1|^2 + \cdots + |a_n|^2) = s_1^2.
    \]
    这说明 $\lVert T \rVert \leqslant s_1$.

    另一方面，取 $v = e_1$，则
    \[
        \lVert e_1 \rVert = 1,\qquad \lVert Te_1 \rVert = \lVert s_1f_1 \rVert = s_1.
    \]
    因而 $\lVert T \rVert \geqslant s_1$. 综上
    \[
        \lVert T \rVert = s_1.
    \]
\end{proof}

也就是说，最大的奇异值就是线性映射对单位向量所能产生的最大伸长倍数. 从几何上看，它对应于前面椭球图像中最长半轴的长度.

\begin{summary}
    奇异值分解的核心思想，是把一般线性映射转化为与正算子 $T^*T$ 有关的问题. 由 $T^*T$ 的谱分解可以得到一组彼此正交的特殊方向，在这些方向上，线性映射只表现为单纯的伸缩，其伸缩倍数就是奇异值.

    对矩阵而言，奇异值分解写成
    \[
        A = U\Sigma V^*,
    \]
    其中 $U, V$ 保长度，$\Sigma$ 负责沿正交方向拉伸. 因而奇异值分解不仅是一个代数分解，更是一个几何分解：单位球在 $A$ 下的像是一个椭球，而奇异值正是这个椭球的半轴长度.

    在应用上，奇异值分解既能帮助我们理解线性变换的几何本质，也能用来刻画算子范数，还能服务于数据降维、图像处理和数值计算等问题. 特征人脸算法便是其中一个典型例子.
\end{summary}

\begin{exercise}
    \exquote[H. Weyl]{在线性代数中，最重要的不是计算本身，而是结构。}

    \begin{exgroup}
        \item 设 $T \in \mathcal{L}(V, W)$. 证明：$T$ 与 $T^*$ 具有完全相同的正奇异值，且这些正奇异值的重数也相同.

        \item 设 $A = \begin{pmatrix}
            1 & 1 \\
            1 & -1
        \end{pmatrix}$. 求 $A$ 的奇异值分解.

        \item 设 $A$ 为 $n$ 阶实矩阵. 证明：$A$ 可逆当且仅当 $A$ 的全部奇异值都非零.
    \end{exgroup}

    \begin{exgroup}
        \item 设 $A$ 为 $m \times n$ 阶复矩阵. 证明：$A^*A$ 与 $AA^*$ 有相同的正特征值，且每个正特征值的重数相同.

        \item 设 $A$ 为 $n$ 阶复矩阵，其奇异值为 $s_1, \ldots, s_n$. 证明：
        \[
            |\det A| = s_1s_2\cdots s_n.
        \]

        \item 设 $A$ 为秩为 $1$ 的 $m \times n$ 阶复矩阵. 证明：存在单位向量 $u \in \mathbf{C}^m,\enspace v \in \mathbf{C}^n$ 和正数 $s$，使得
            \[
                A = suv^*.
            \]
    \end{exgroup}

    \begin{exgroup}
        \item 设 $A \in \mathbf{R}^{m \times n}$ 的奇异值为 $s_1 \geqslant \cdots \geqslant s_r > 0$. 证明：单位球
        \[
            B = \{x \in \mathbf{R}^n \mid \lVert x \rVert \leqslant 1\}
        \]
        在 $A$ 下的像是一个椭球，并指出该椭球各条半轴的方向与长度.

        \item 设 $A$ 为 $n$ 阶正规矩阵，特征值为 $\lambda_1, \ldots, \lambda_n$. 证明：$A$ 的奇异值恰好是
        \[
            |\lambda_1|, |\lambda_2|, \ldots, |\lambda_n|.
        \]

        \item 设 $A = U\Sigma V^*$ 是矩阵 $A$ 的奇异值分解，$\Sigma$ 的对角线上前 $r$ 个元素为正，其余为 $0$. 证明：
        \[
            A = s_1u_1v_1^* + s_2u_2v_2^* + \cdots + s_ru_rv_r^*,
        \]
        其中 $u_i$ 与 $v_i$ 分别是 $U$ 与 $V$ 的第 $i$ 列.
    \end{exgroup}
\end{exercise}
